%=====================================================================
% 10  费米子的量子力学
% 原文：Section 10. Quantum mechanics for fermions
%       （arXiv:2409.03117v1, p. 133--141）
%=====================================================================
\section{费米子的量子力学}

\begin{keytakeaways}
\begin{itemize}
  \item \textbf{物理图像：}把 Feynman 演算推广到超情形，偶边与奇边分别对应玻色与费米自由度；当作用量对费米变量为二次型时，奇变量可精确积掉，得到含 Pfaffian 的纯玻色积分，此时奇线组成费米圈。
  \item \textbf{数学工具：}超向量空间与超迹、Grassmann 变量的 Berezin 积分、Clifford 代数及其不可约表示（偶维时唯一；奇维时分解为两个不可约分量，但作为超表示仍不可约）。
  \item \textbf{关键结论：}费米量子力学的经典状状态空间不含动量变量，Poisson 括号由内积空间上的对称型给出，量子化得到 Clifford 代数；哈密顿量的本征谱给出费米振子的能谱，环面配分函数由超迹给出。
  \item \textbf{易错点：}奇变量的交换带来符号；Wick 转动后虚单位不消失；奇维时配分函数为零；每个费米圈贡献一个负号。
\end{itemize}
\end{keytakeaways}

\bigskip
\noindent\rule{\textwidth}{0.6pt}
\bigskip

\subsection{超情形下的 Feynman 演算}

Wick 定理使我们能把 Feynman 演算推广到超情形。具体地说，设
$$
V=V_0\oplus V_1
$$
是带超体积元 $dv=dv_0(dv_1)^{-1}$ 的有限维实超空间，配有对称非退化形式 $B=B_0\oplus B_1$（$B_0>0$）。设
$$
S(v)=\frac{1}{2}B(v,v)-\sum_{r\ge 3}\frac{B_r(v,v,...,v)}{r!}
$$
是 $V$ 上的偶函数（作用量）。注意 $B_r$（$r\ge 3$）可以含有同时涉及奇变量与偶变量的混合项，例如 $x\xi_1\xi_2$（即所谓的“Yukawa 项”）。我们考虑积分
$$
I(\hbar)=\int_V \ell_1(v_0)...\ell_n(v_0)\lambda_1(v_1)...\lambda_p(v_1)e^{-\frac{S(v)}{\hbar}}dv,
$$
其中 $v_0,v_1$ 是 $v$ 的偶分量与奇分量。于是该积分有按 Feynman 图写出的 $\hbar$ 展开。由于 $v$ 同时含有奇部与偶部，这些图将含有“奇”边与“偶”边（通常分别用直线与波浪线表示）。更精确地说，我们写
$$
B_r(v,v,...,v)=\sum_{s=0}^r\begin{pmatrix}r\\ s
\end{pmatrix}
B_{s,r-s}(v_1,...,v_1,v_0,...,v_0),
$$
其中 $B_{s,r-s}$ 关于 $v_1$ 的齐次次数为 $s$，关于 $v_0$ 的齐次次数为 $r-s$（亦即仅当 $s$ 为偶数时它才非零）。于是对每一项 $B_{s,r-s}$，我们指定一朵 $(s,r-s)$ 价花，即一朵带 $s$ 条奇出边与 $r-s$ 条偶出边的花；并且对奇出边的集合指定哪些排序为偶。然后，给定花的一种排布，对出边的每个配对 $\sigma$，我们可以通过缩并张量 $B_{s,r-s}$ 来定义振幅 $\Bbb F(\sigma)$（并小心处理符号）。容易验证，给出同一图的所有配对将以相同符号贡献于 $I(\hbar)$，因此我们得到与玻色情形几乎相同的公式：
$$
I(\hbar)=(2\pi)^{\frac{\dim V_0}{2}}\hbar^{\frac{\dim V_0-\dim V_1}{2}}
\frac{{\rm Pf}(-B_1)}{\sqrt{\det B_0}}
\sum_{\Gamma}\frac{\hbar^{b(\Gamma)}}{|{\rm Aut}(\Gamma)|}\Bbb F_\Gamma (\ell_1,...,\ell_n,\lambda_1,...,\lambda_p),
$$
其中求和取遍带 $n$ 条偶出边与 $p$ 条奇出边的图。

\begin{remark} 更精确地说，配对 $\sigma$ 的符号 $\varepsilon_\sigma$ 可如下定义：把出边从第一朵花、第二朵花等等开始依次标号为 $1,2,...$，使得每朵花上的标号均为偶；然后把标号写成一个序列，按任意顺序枚举由 $\sigma$ 定义的配对（两个标号中较小者在前）。符号 $\varepsilon_\sigma$ 按定义就是这个排序（作为 $1,2,...$ 的置换）的符号。于是 $\Bbb F_\Gamma$ 取为任一给出 $\Gamma$ 且为\emph{正}的配对 $\sigma$ 的 $\Bbb F(\sigma)$，即满足 $\varepsilon_\sigma=1$；对于负配对，$\Bbb F_\Gamma=-\Bbb F(\sigma)$。
\end{remark}

% [Tutor Note: 这里的符号约定值得核对：对一条含 2k 条边的费米圈（k≥1），沿圈缩并要作 2k−1 次相邻奇变量的交换，故该圈贡献 (−1)^{2k−1}=−1，这正是“偶循环是奇置换”的实质。注意 (−1)^r 中的 r 是费米圈总数；若改变奇出边的定向或缩并次序，整体符号可能改变，但配分函数与关联函数在等价重定义下不变，核对时应把所有圈的因子一并计入。]

在物理学中考虑的大多数（但并非全部）情形，作用量对费米变量是二次的，即
$$
S(v)=S_b(v_0)-\tfrac{1}{2}S_f(v_0)(v_1,v_1),
$$
其中 $S_f(v_0)$ 是 $\Pi V_1$ 上的反对称双线性形式。此时利用费米 Wick 定理，我们可以精确地对 $v_1$ 积分，并把 $I(\hbar)$ 化为纯玻色积分。例如，若只有 $\ell_i$ 而没有 $\lambda_i$，则
$$
I(\hbar)=\hbar^{-\frac{\dim V_1}{2}}\int_{V_0}\ell_1(v_0)...\ell_n(v_0)e^{-\frac{S_b(v_0)}{\hbar}}
{\rm Pf}(S_f(v_0))dv_0.
$$
在这种情形，所有带奇出边的顶点都恰有两条奇出边，因此在任何只有偶出边的 Feynman 图中，奇线构成互不相交的简单曲线，称为\emph{费米圈}（事实上，最后一个公式不过是把这些圈视为一种新顶点的结果——请自行验证！）。此时有如下给 Feynman 图指定符号的简单办法：对每个带两条奇出边的顶点，把第一条边定向为指向内、第二条边定向为指向外；我们只允许保持定向的连接（配对）（于是费米圈成为有向的）。则符号为 $(-1)^r$，其中 $r$ 是费米圈的数目（亦即每个费米圈贡献一个负号）。这可由“偶循环是奇置换”这一事实推出。

\subsection{费米量子力学} \label{fqm}

现在从有限维费米积分过渡到量子力学，即关于一个（偶的）实变量 $t$ 的费米函数的积分。

先讨论拉格朗日框架下的费米经典力学。它与玻色情形的差别在于：粒子的“轨迹”由一个单变量的\emph{奇值}函数描述，即 $\psi: \Bbb R\to \Pi V$，其中 $V$ 是向量空间。数学上这意味着场（=轨迹）的空间是奇向量空间 $\Pi C^\infty(\Bbb R,V)$。拉格朗日量 ${\mathcal L}(\psi)$ 是这种场的一个局部表达式（即 $\psi,\dot{\psi},...$ 的多项式），作用量是积分 $S=\int_{\Bbb R}{\mathcal L}dt$。这意味着作用量是空间 $\Lambda (C^\infty_0(\Bbb R,V)^*)$ 中的一个元素。

例如考虑单个标量值自由费米子 $\psi(t)$ 的理论。按定义，这种理论的拉格朗日量为
$$
{\mathcal L}=\frac{1}{2}
\psi\dot{\psi},
$$
亦即作用量为
$$
S=
\frac{1}{2}\int \psi\dot{\psi} dt.
$$
该拉格朗日量是自由粒子拉格朗日量 $\frac{\dot{q}^2}{2}$ 的奇类比。

\begin{remark} 注意 $\psi\dot{\psi}\ne \frac{d}{dt}(\frac{\psi^2}{2})=0$，因为 $\psi\dot{\psi}=-\dot{\psi}\psi$，所以这个拉格朗日量是“合理的”。另一方面，同一个拉格朗日量在玻色情形却不合理，因为它会是全导数，从而作用量为零。最后注意，在费米情形使用通常的玻色拉格朗日量 $\frac{1}{2}(\dot{q}^2-m^2q^2)$ 同样不合理；若 $q$ 取奇值，它将恒为零。
\end{remark}

拉格朗日量 ${\mathcal L}$ 在重参数化群 ${\rm Diff}_+(\Bbb R)$ 下不变，且该拉格朗日量的欧拉–拉格朗日方程为
$$
\dot{\psi}=0
$$
（即没有动力学）。具有这种性质的理论称为\emph{拓扑量子场论}。

% [Tutor Note: 奇变量的欧拉–拉格朗日方程须用左导数（或与之等价的统一符号约定）推得：由 L=(1/2)ψψ̇ 变分得 δS=∫δψ ψ̇，故方程是 ψ̇=0；若照搬玻色情形的右导数并随意交换奇变量的次序，容易多出整体负号。一般的规则是：含奇变量的变分与 Berezin 积分都按“左作用”约定进行，所有交换各计一次 (−1)。]

现在转向拉格朗日框架下的量子理论，即由 Feynman 积分 $\int \psi(t_1)...\psi(t_n)e^{\frac{iS(\psi)}{\hbar}}D\psi$ 给出的理论。在玻色情形，我们把这类表达式在空间 $C^\infty_0(\Bbb R)$ 上“积分”。由于无限维测度论的困难，这种积分并没有直接意义，因此我们必须借助 Wick 公式与 Feynman 图，用 $\hbar$ 展开来赋予其意义。在费米情形，情况类似。具体地说，现在必须在 $\Pi C^\infty_0(\Bbb R)$ 上积分函数，这些函数是 $\Lambda {\mathcal D}(\Bbb R)$ 的元素，其中 ${\mathcal D}(\Bbb R)$ 是 $\Bbb R$ 上分布的空间。尽管费米情形不需要测度论（因为积分完全是代数的），我们仍有定义积分的困难：回忆按定义积分应是被积函数作为 $\Lambda {\mathcal D}(\Bbb R)$ 中元素的最高次分量，而这没有意义，因为无限维空间的外代数没有最高次分量。因此我们必须采用与玻色情形相同的策略，即 Feynman 图。

例如，让我们定义标量值自由费米子的量子理论，即由拉格朗日量 ${\mathcal L}=\frac{1}{2}\psi\dot{\psi}$ 描述的理论。按照我们在玻色情形所用的要诀，该理论的两点函数 $\langle\psi(t_1)\psi(t_2)\rangle$ 应是函数 $G(t_1-t_2)$，其中 $G$ 是微分方程
$$
\frac{dG}{dt}=i\delta(t).
$$
的解（因子 $i$ 来自 Feynman 积分中的指数；注意在费米情形，它在 Wick 转动下\emph{不会}消失）。

该方程的通解具有形式
$$
G(t)=\frac{i}{2}{\rm sign}(t)+C.
$$
由于场 $\psi(t)$ 的费米性质，自然要求 $G(-t)=-G(t)$，即关联函数是反对称的；这唯一确定了解 $G(t)=\frac{i}{2}{\rm sign}(t)$（由该条件我们也看出应令 $G(0)=0$）。通常，$2n$ 点关联函数由 Wick 公式定义。也就是说，对互异的 $t_j$，
$$
\langle\psi(t_1)...\psi(t_{2n})\rangle=(-1)^\sigma(2n-1)!!(\tfrac{i}{2})^n,
$$
其中 $\sigma$ 是把 $t_j$ 按降序排列的置换。若至少有两个点重合，则关联函数为零。

于是我们看到，关联函数在 ${\rm Diff}_+(\Bbb R)$ 下不变。换言之，用物理术语说，我们得到一个\emph{拓扑量子场论}。

注意该模型在欧氏情形下的关联函数与 Minkowski 情形相同，因为它们是关于 $t_j$ 的（分段）常数。特别地，它们在无穷远处不衰减，因此我们的理论不具有聚类性质。

% [Tutor Note: 这里 G(−t)=−G(t) 的反对称性来自费米场的奇性，它同时固定 G(0)=0 并排除解中的常数 C。注意 Wick 转动 t→−iτ 时因子 i 不会消失（玻色情形则会给出指数衰减的 e^{−mτ}），因此欧氏关联函数仍是分段常数而非衰减的；这正是无质量费米理论没有聚类性质的表现。]

我们考虑了无质量费米场论。现在考虑有质量情形。这意味着要在拉格朗日量中加入一个不含导数的 $\psi$ 的二次项。若只有一个场 $\psi$，这样的项唯一的选择是 $\psi^2$，而它为零。因此在有质量情形必须至少有两个场。于是考虑两个费米场 $\psi_1,\psi_2$ 的理论，其（欧氏）拉格朗日量为
$$
{\mathcal
L}=\frac{1}{2}(\psi_1\dot{\psi}_1+\psi_2\dot{\psi_2}-m\psi_1\psi_2),
$$
其中 $m>0$ 是质量参数。该模型的 Green 函数满足微分方程
$$
\frac{dG}{dt}-MG=i\delta(t),
$$
其中 $M=\begin{pmatrix} 0&m\\ -m&0\end{pmatrix}$，$G$ 是二阶矩阵值函数。该方程的通解为
$$
G(t)=\begin{cases}e^{Mt}Q_-,\ t<0\\ e^{Mt}Q_+,\ t>0\end{cases}
$$
其中 $Q_+-Q_-=i$。现在，我们希望 Wick 转动后的 Green 函数 $G(-it)$ 具有聚类性质。因此要求
$$
\lim_{t\to +\infty}e^{-iMt}Q_+=0,\ \lim_{t\to -\infty}e^{-iMt}Q_-=0.
$$
这蕴含 $Q_+=iP_+$，$Q_-=-iP_-$，其中 $P_{\pm}$ 是到 $iM$ 的具有本征值 $\pm m$ 的本征空间的正交投影算子（并且 $G(0)=0$）。

\begin{remark} 容易把这一分析推广到如下情形：$\psi$ 取值于正定内积空间 $V$，而 $M:V\to V$ 是反对称算子，因为这种情形是上述情形的直和。
\end{remark}

当 $M$ 非退化时，可以定义相应的带相互作用的理论，即含有高于二次的 $\psi$ 项的理论。具体地说，把关联函数定义为相应 Feynman 图振幅之和。与上面讨论过的有限维情形类比，我们把这一工作留给读者。

\subsection{超 Hilbert 空间}

量子系统的状状态空间是 Hilbert 空间。由于我们打算对费米子做哈密顿量子力学，必须定义这一概念的超类比。

设 ${\mathcal H}={\mathcal H}_0\oplus {\mathcal H}_1$ 是 $\Bbb Z/2$-分次复向量空间。

\begin{definition} (i) ${\mathcal H}$ 上的 Hermitian 形式是一个偶的半双线性形式 $\la,\ra$，使得对偶的 $x,y$ 有 $\la x,y\ra=\overline{\la y,x\ra}$，而对奇的 $x,y$ 有 $\la x,y\ra=-\overline{\la y,x\ra}$。

(ii) 若对偶的 $x\ne 0$ 有 $\la x,x\ra>0$，且对奇的 $x\ne 0$ 有 $-i\la x,x\ra>0$，则称 Hermitian 形式是正定的。超 Hilbert 空间是配有正定 Hermitian 形式 $\la,\ra$ 的超空间，且它在相应的范数下完备。

(iii) 设 $\mathcal H$ 是超 Hilbert 空间，$T: {\mathcal H_0}\oplus \Pi {\mathcal H_1}\to 
{\mathcal H_0}\oplus \Pi {\mathcal H_1}$ 是底层的纯偶空间之间的齐次线性算子。Hermitian 伴随算子 $T^\dagger$ 由方程 $\la x,T^\dagger y\ra=(-1)^{p(x)p(T)}\la Tx,y\ra$ 定义，其中 $p$ 表示宇称。
\end{definition}

\subsection{费米量子力学的哈密顿形式}

现在讨论自由费米子理论的哈密顿图像应当是什么。更精确地说，设 $V$ 是正定有限维实内积空间，考虑拉格朗日量
$$
{\mathcal L}=\frac{1}{2}((\psi,\dot{\psi})-(\psi,M\psi)),
$$
其中 $\psi: \Bbb R\to \Pi V$，而 $M:V\to V$ 是反对称算子。

为理解哈密顿图像应当是什么，我们与玻色情形作比较。具体地说，考虑拉格朗日量
$$
{\mathcal L}_b=\frac{1}{2}(\dot{q}^2-m^2q^2),
$$
其中 $q:\Bbb R\to V$。此时经典状状态空间是
$$
Y:=T^*V=V\oplus
V^*.
$$
运动方程是 Newton 方程
$$
\ddot{q}=-m^2q,
$$
它们可化为 哈密顿方程
$$
\dot{q}=p,\ \dot{p}=-m^2q.
$$
经典可观测量的代数是 $C^\infty(Y)$，其 Poisson 括号定义为 $\lbrace{a,b\rbrace}=(a,b)$（$a,b\in Y^*$），其中 $(,)$ 是 $Y^*$ 上与 $Y$ 上自然辛形式互逆的形式。哈密顿量 $H$ 由“运动方程为 $\dot{f}=\lbrace{f,H\rbrace}$”这一条件确定（至多差一个常数）；在此情形它等于 $H=\frac{1}{2}(p^2+m^2q^2)$。

费米情形类似，但有一些重要的差别，我们下面说明。具体地说，容易算出运动方程（即欧拉–拉格朗日方程）为
$$
\dot{\psi}=M\psi.
$$
与玻色情形的主要差别在于：该方程是一阶而非二阶的，因此经典状状态空间就是 $\Pi V$（不引入动量或速度变量）。于是经典可观测量的代数是 $C^\infty(\Pi V)=\Lambda V^*$。为在该代数上定义 Poisson 括号，回忆 $\Pi V$ 具有由 $V$ 上的\emph{对称}形式 $(,)$ 定义的自然“辛结构”。于是可用与上面相同的公式在 $\Lambda V^*$ 上定义 Poisson 括号：当 $a,b\in V^*$ 时 $\lbrace{a,b\rbrace}=(a,b)$。更精确地说，$\lbrace{,\rbrace}$ 是 $\Lambda V^*$ 上唯一的（超意义下的）反对称双线性运算，它对 $a,b\in V^*$ 限制为 $(a,b)$，且对每个变量都是导子：
$$
\lbrace{a,bc\rbrace}=\lbrace{a,b\rbrace}c+(-1)^{p(a)p(b)}b\lbrace{a,c\rbrace},
$$
其中 $p(a)$ 表示 $a$ 的宇称。

% [Tutor Note: 费米情形所谓的“辛结构”由 V 上的对称形式给出，初看似乎与“辛=反对称”矛盾。原因在于奇坐标本身反对易：若 ξ^i 是对应于 V 的奇坐标，则超 Poisson 括号的反对称性是 {a,b}=−(−1)^{p(a)p(b)}{b,a}，于是对称的 {ξ^i,ξ^j}=(e_i,e_j) 与超反对称并不冲突。展开 ΛV* 时所有奇变量两两交换产生的 (−1) 正是这样来的。]

现在容易看出哈密顿量应当扮演什么角色。更精确地说，用 Legendre 变换的定义在我们的情形不适用，因为 Legendre 变换是对速度变量作的，而费米情形没有速度变量。另一方面，正如我们在第 8 章讨论过的，在玻色情形运动方程
$$
\dot{f}=\lbrace{f,H\rbrace}
$$
唯一确定 $H$（至多差一个常数）。费米情形也是如此。具体地说，观察运动方程 $\dot{\psi}=M\psi$，容易看出哈密顿量等于
$$
H=\frac{1}{2}(\psi,M\psi).
$$
特别地，若 $M=0$（无质量情形），则哈密顿量为零（这是拓扑场论的一个特征性质）。

现在转向量子理论。在玻色情形，量子可观测量的代数是 $C^\infty(Y)$ 的非交换形变，其中关系 $\lbrace{a,b\rbrace}=(a,b)$ 被替换为其量子类比
$$
ab-ba=i(a,b)
$$
（至多差 普朗克常数因子，这里我们令其为 $1$）。特别地，多项式可观测量的子代数是 Weyl 代数 $W(Y)$，它由 $Y^*$ 以此定义关系生成。与此类比，我们应把费米情形的量子可观测量代数定义为由 $V^*$ 以关系
$$
ab+ba=i(a,b)
$$
生成的代数（它形变定义了 $\Lambda V^*$ 的关系 $ab+ba=0$）。于是我们回顾如下定义。

\begin{definition} 设 $V$ 是域 $k$ 上的向量空间，配有对称双线性形式 $Q$。\emph{Clifford 代数} ${\rm Cl}(V,Q)$ 是由 $V$ 以定义关系 $ab+ba=Q(a,b)$（$a,b\in V$）生成的代数。
\end{definition}

我们看到，量子可观测量代数应是 ${\rm Cl}(V_{\Bbb C}^*,i(,))$。注意与经典情形一样，该代数自然地是 $\Bbb Z/2$ 分次的，因此我们有偶的与奇的量子可观测量。

现在来看量子态对应的 Hilbert 空间应当是什么。在玻色情形它是 $L^2(V)$，而由著名的 Stone–von Neumann 定理，它是 $W(Y)$ 的唯一不可约酉表示。与此类比，在费米情形，状态 Hilbert 空间应是 ${\rm Cl}(V^*_{\Bbb C})$ 在超向量空间 ${\mathcal H}$ 上的不可约酉表示。

Clifford 代数 ${\rm Cl}(V^*_{\Bbb C})$ 的结构是众所周知的。具体地说，分别考虑 $\dim V$ 为奇数与偶数的情形。

偶情形：$\dim V=2d$，${\rm Cl}(V^*_{\Bbb C})$ 是单代数（即同构于矩阵代数），并有唯一的不可约表示 ${\mathcal H}$，维数为 $2^d$。该表示构造如下：选取分解 $V_{\Bbb C}=L\oplus L^*$，其中 $L,L^*$ 是拉格朗日子空间；则 ${\mathcal H}=\Lambda L$，其中 $L\subset V^*_\Bbb C$ 通过乘法作用，而 $L^*$ 通过微分作用（乘以 $-i$）。${\mathcal H}$ 上的超空间结构就是外代数上的标准结构。

奇情形：$\dim V=2d+1$，选取分解
$$
V_{\Bbb C}=
L\oplus L^*\oplus K,
$$
其中 $L,L^*$ 是极大迷向子空间，$K$ 是与 $L$ 和 $L^*$ 正交的非退化一维子空间。令 ${\mathcal H}=\Lambda(L\oplus K)$，其中 $L,K$ 通过乘法作用，$L^*$ 通过（$-i$ 倍）微分作用。这是 ${\rm Cl}(V^*_{\Bbb C})$ 的一个带 ${\Bbb Z/2}$ 分次的表示。该表示不是不可约的，它分解为两个不同构不可约表示的直和 ${\mathcal H}_+\oplus {\mathcal
H}_-$（这与如下事实有关：奇数 $\dim V$ 时的 Clifford 代数不是单代数，而是两个单代数（即矩阵代数）的直积）。然而，该分解与 ${\Bbb Z/2}$ 分次不相容，因此作为超表示，${\mathcal H}$ 是不可约的。

现在容易证明：无论奇情形还是偶情形，空间 ${\mathcal H}$ 都带有唯一的（至多差一个缩放）Hermitian 形式，使得 $V^*\subset V^*_{\Bbb C}$ 通过自伴算子作用。该形式是正定的。因此对任意 $\dim V$，情形都与玻色情形类似。

现在来看 ${\mathcal H}$ 上哪个算子应扮演系统的哈密顿量。最自然的选择是把量子哈密顿量定义为经典哈密顿量 $H=\frac{1}{2}(\psi,M\psi)$ 的显然量子化。具体地说，若 $\varepsilon_i$ 是 $V^*$ 的正交基，$a_{ij}$ 是 $M$ 在该基下的矩阵，则令
$$
\widehat{H}=\frac{1}{2}\sum_{i,j} a_{ij}\varepsilon_i\varepsilon_j.
$$
为更显式地计算该算子，我们（不失一般性）假定所选取的 $V_{\Bbb C}$ 的分解在 $M$ 下稳定。设 $\xi_j$ 是 $M$ 在 $L$ 中的本征基，本征值为 $im_j$（$m_j\ge 0$），而 $\partial_j$ 是沿该基向量的微分。则
$$
\widehat H=\sum_j m_j(\xi_j\partial_j-\partial_j\xi_j)=
\sum_j m_j(2\xi_j\partial_j-1).
$$

这表明，若 $\dim V$ 为偶数，则我们的理论在长度为 $L$ 的圆上的配分函数为
$$
Z={\rm sTr}(e^{-L\widehat H})=\prod_j (e^{m_jL}-e^{-m_jL}).
$$
若 $V$ 的维数为奇数，则配分函数为零。

% [Tutor Note: 超迹须按 sTr(A)=Tr((−1)^F A) 计算：对每个费米振子，Ĥ 在占据数为 0、1 的态上的本征值分别为 −m_j、+m_j，故 e^{−LĤ} 的贡献为 e^{m_jL} 与 −e^{−m_jL}，合计 e^{m_jL}−e^{−m_jL}。这正是费米振子与玻色振子配分函数不同之处（后者给 1/(2sinh(m_jL/2)) 型因子）。奇维时超迹为零，反映 Ĥ 的谱在偶部与奇部成对出现，即未破缺的超对称配对。]

现在我们想考虑 Feynman–Kac 公式的费米类比。为简单起见，考虑全有质量的情形，即 $\dim V$ 为偶数且 $m_j\ne 0$（亦即 $M$ 非退化）。在此情形，$\widehat H$ 有唯一的（至多差一个缩放）最低本征向量，即 $\Omega=1$。

设 $\psi(0)\in V\otimes {\rm End}({\mathcal H})$ 是对应于作用映射 $V^*\to {\rm End}({\mathcal H})$ 的元素（即\emph{Clifford 乘法}），并令 $\psi(t)=e^{it\widehat H}\psi(0)e^{-it\widehat H}$。又用 $\langle\psi(t_1)...\psi(t_n)\rangle$（$t_1\ge...\ge t_n$）表示拉格朗日框架下自由理论的关联函数，取值于 $V^{\otimes n}$（于是在该表达式中 $\psi(t_j)$ 只是形式符号而非算子）。

\begin{theorem} (Feynman-Kac 公式)
(i) 对直线上的自由理论，我们有
$$
\langle\psi(t_1)...\psi(t_n)\rangle=
\la\Omega,\psi(t_1)...\psi(t_n)\Omega\ra.
$$

(ii) 对长度为 $L$ 的圆上的自由理论，我们有
$$
\langle\psi(t_1)...\psi(t_n)\rangle=
\frac{{\rm sTr}(\psi(t_1)...\psi(t_n)e^{-L\widehat H})}{{\rm sTr}(e^{-L\widehat H})}.
$$
\end{theorem}

\begin{exercise} 证明该定理。（证明与自由情形的定理 \ref{FK} 类似。）
\end{exercise}

% [Tutor Note: 该定理是第 8 章 Feynman–Kac 公式的超类比：把“对基态取期望”替换为“除以配分函数的超迹”。要点有二：(1) 由于 Ĥ 在费米振子的奇态上反号，e^{−LĤ} 的符号必须由超迹中的 (−1)^F 补偿，否则环面上的关联函数不满足平移不变性；(2) 直线情形要求 t_1≥...≥t_n 并取真空期望，因为费米场算符在编时乘积中交换时要产生符号。自由情形的证明可用费米高斯积分（Pfaffian/Berezin 积分）逐项比较完成。]

现在读者应当能直接叙述并证明包含玻色与费米有质量场的相互作用（即不一定自由）量子力学模型的 Feynman–Kac 公式。我们把它留作一个有启发性的习题。

\begin{exercise} (i) 考虑带 Yukawa 耦合的量子力学。即有一个标量玻色子 $\phi(t)$ 与两个费米子 $\psi_1(t),\psi_2(t)$，其欧氏拉格朗日量为
$$
{\mathcal L}=\frac{1}{2}(\dot{\phi}^2+m^2\phi^2+
\psi_1\dot{\psi_1}+\psi_2\dot{\psi_2}-\mu \psi_1\psi_2)+
g\phi\psi_1\psi_2.
$$
计算两点函数 $\langle\phi(t)\phi(0)\rangle$ 模 $g^3$（在欧氏情形下）。

\textbf{提示.} 对自由理论答案的修正由一张 Feynman 图给出。不要忘记自同构群以及与费米圈相应的负号。

(ii) 在同一理论中，计算两点函数 $\langle\psi_1(t)\psi_1(0)\rangle$ 模 $g^3$（在欧氏情形下）。相应的图是否有非平凡自同构？
\end{exercise}
